%!TEX root = ../../../exa-ma-d7.1.tex

\section{Double Mach reflection on an adaptive mesh}
\label{sec:app:specs:app-samurai-euler-dmr}

This mini-app measures the strong scalability of samurai (see
\Cref{sec:Samurai:software}), and its throughput when the resolution grows
with the number of processes, on the double Mach reflection, a standard test
case for the compressible Euler equations. Unlike the Burgers benchmark of
\Cref{sec:app:specs:app-samurai-burgers-2d}, the adapted mesh is not balanced
a priori: the refined cells gather along the shocks, the contact
discontinuity and the jet near the wall, and their distribution changes as the
flow develops. The mini-app is therefore also the test case on which we compare
the load balancing strategies of samurai.

\Cref{tab:app-samurai-euler-dmr} describes the specifications of the
application.

\begin{table}[ht]
    \centering
    \begin{tblr}{
        colspec = {l X[10cm]},
        row{odd} = {numpexlightergray},
        hlines = {0.1pt, numpexgray},
        vlines = {numpexgray},
        row{1} = {numpexgray, fg=white, font=\bfseries},
    }
        Field & Details \\
        id & \texttt{app-samurai-euler-dmr} \\
        name & Double Mach reflection, strong scaling and scaled-resolution series \\
        Partners & IPP \\
        PC & PC1 - ExaMA \\
        Responsible (Permanent) & L. Gouarin \\
        WP7 Engineer &  \\
        work\_package & WP1 \\
        application\_type & mini-app \\
        purpose & Strong scalability, and throughput at a resolution growing with the number of processes, of an adaptive finite volume solver for the compressible Euler equations with strong shocks, and comparison of the dynamic load balancing strategies of samurai \\
        Method-Algorithm WP1 & mesh adaptation, multiresolution analysis, finite volume, dynamic load balancing \\
        Method-Algorithm WP2 & \\
        Method-Algorithm WP3 & \\
        Method-Algorithm WP4 & \\
        Method-Algorithm WP5 & \\
        Method-Algorithm WP6 & \\
        WP7 & \\
        metrics & \SetCell{l} \texttt{strong-scalability}, \texttt{time-to-solution}, \texttt{regression-test-pass} \\
        status & benchmark-in-progress \\
        Benchmark scope & multi-node \\
        Framework & samurai \\
        parallel\_framework & MPI \\
        repo\_url & \url{https://github.com/hpc-maths/samurai-euler}\\
    \end{tblr}
    \caption{Description of the mini-app \texttt{app-samurai-euler-dmr}.}
    \label{tab:app-samurai-euler-dmr}
\end{table}


%%%%%%%%%%%%%%%%%%%%%%%%%%%%%%%%%%%%%%%%%%%%%%%%%%%%%%%%%%%%%%%%%%%%%%%%%%%%%%%%%%%%%%%%%%%%%%%%%%%%%%%%%%%%%%%%%%%%%%%%


\subsection{Description of the benchmark}

We solve the two-dimensional compressible Euler equations
\begin{equation}\label{eq:specs:samurai:euler}
  \partial_t
  \begin{pmatrix} \rho \\ \rho \mathbf{u} \\ \rho E \end{pmatrix}
  + \nabla\cdot
  \begin{pmatrix} \rho \mathbf{u} \\ \rho \mathbf{u}\otimes\mathbf{u} + p\,\mathbf{I} \\ (\rho E + p)\,\mathbf{u} \end{pmatrix}
  = 0,
  \qquad
  p = (\gamma - 1)\left(\rho E - \frac{1}{2}\rho |\mathbf{u}|^2\right),
\end{equation}
for an ideal gas with $\gamma = 1.4$. The set-up is the one of section IVa of
\citeA{woodward_numerical_1984}. A Mach 10 shock hits a
30 degree ramp, represented by a reflecting wall along the bottom of the domain
$[0,4]\times[0,1]$ from $x_0 = 1/6$. At $t = 0$, the shock makes an angle of
$60^\circ$ with the $x$-axis and passes through $(x_0, 0)$. The gas at rest in
front of it is in the state $\mathbf{U}_R$, with $(\rho, p) = (1.4, 1)$, so
that its sound speed is 1. The Rankine-Hugoniot conditions give the post-shock
state $\mathbf{U}_L$, with $\rho = 8$, $p = 116.5$ and a velocity of magnitude
$8.25$ normal to the shock.

\Cref{fig:specs:samurai:dmr:setup} shows the initial condition and the boundary
conditions. The left boundary and the part of the bottom boundary upstream of
$x_0$ impose $\mathbf{U}_L$. Downstream of $x_0$, the bottom boundary is a
reflecting wall. The top boundary follows the exact position of the shock,
\begin{equation}\label{eq:specs:samurai:dmr:shock}
  x_s(t) = x_0 + \frac{1 + 20\,t}{\sqrt{3}},
\end{equation}
and imposes $\mathbf{U}_L$ for $x < x_s(t)$ and $\mathbf{U}_R$ beyond. The right
boundary is an outflow.

\begin{figure}[ht]
  \centering
  \begin{tikzpicture}[x=3.1cm, y=3.1cm, font=\small]
    \def\xz{0.16667}
    \def\xs{0.74402}
    % initial states on each side of the shock
    \fill[customorange!25] (0,0) -- (\xz,0) -- (\xs,1) -- (0,1) -- cycle;
    \fill[customdarkblue!8] (\xz,0) -- (4,0) -- (4,1) -- (\xs,1) -- cycle;
    \node at (0.25,0.78) {$\mathbf{U}_L$};
    \node at (2.3,0.5) {$\mathbf{U}_R$};
    % initial shock, angle with the x-axis and post-shock velocity
    \draw[very thick, customorange!80!black] (\xz,0) -- (\xs,1);
    \draw (\xz+0.16,0) arc (0:60:0.16);
    \node at (\xz+0.24,0.11) {\footnotesize $60^\circ$};
    \draw[->, thick] (0.10,0.50) -- ++(0.866*0.22,-0.5*0.22) node[below left=-1pt] {\footnotesize $\mathbf{u}_L$};
    % domain
    \draw[thick] (0,0) rectangle (4,1);
    % bottom: imposed state upstream of x0, reflecting wall downstream
    \draw[ultra thick, customorange!80!black] (0,0) -- (\xz,0);
    \draw[ultra thick] (\xz,0) -- (4,0);
    \foreach \x in {0.25,0.35,...,4.0} \draw (\x,0) -- ++(-0.05,-0.06);
    \node[below left=2pt and -1pt] at (\xz,0) {$\mathbf{U}_L$};
    \draw (\xz,0) -- (\xz,-0.2) node[below] {$x_0$};
    \node[below=6pt] at (2.1,0) {reflecting wall};
    % top: exact shock position x_s(t)
    \draw (\xs,1) -- (\xs,1.1) node[above] {$x_s(t)$};
    \draw[->] (\xs+0.03,1.05) -- (\xs+0.35,1.05);
    \node[above=2pt] at (\xs/2,1) {$\mathbf{U}_L$};
    \node[above=2pt] at (2.3,1) {$\mathbf{U}_R$};
    % left and right
    \node[left] at (0,0.5) {$\mathbf{U}_L$};
    \node[right, rotate=90, anchor=north] at (4.04,0.5) {outflow};
  \end{tikzpicture}
  \caption{Double Mach reflection: initial condition and boundary conditions on
    the domain $[0,4]\times[0,1]$. $\mathbf{U}_L$ is the post-shock state, $\rho = 8$, $p = 116.5$,
    $\mathbf{u}_L = 8.25\,(\sin 60^\circ, -\cos 60^\circ)$, and $\mathbf{U}_R$
    the gas at rest, $\rho = 1.4$, $p = 1$, $\mathbf{u}_R = 0$. The shock starts
    at $x_0 = 1/6$ on the bottom boundary, and $x_s(t)$ is its exact position on
    the top boundary, given by \eqref{eq:specs:samurai:dmr:shock}.}
  \label{fig:specs:samurai:dmr:setup}
\end{figure}

The simulation runs until $T_f = 0.2$. The reflected shocks, the Mach stem and
the jet that rolls up along the wall are thin structures in a large domain at
rest or in uniform motion, which is the configuration where mesh adaptation
pays off.

\paragraph{Numerical method.}
The fluxes are computed with the HLLC Riemann solver on face values
reconstructed by MUSCL with the monotonized central limiter, on a stencil of
four cells. The time integration is a Strang splitting,
$X(\Delta t/2)\,Y(\Delta t)\,X(\Delta t/2)$, in which each directional sweep
uses the Hancock predictor, with $\mathrm{CFL} = 0.4$. The mesh is adapted at
every time step by multiresolution analysis
\cite{cohen_fully_2003,bellotti_multidimensional_2022} on relative details,
between the levels $\ell_{\min}$ and $\ell_{\max}$, with the threshold set by
\texttt{-{}-mr-eps}. The mesh adaptation uses a prediction operator specific to the
Euler variables.

\paragraph{Load balancing.}
Since the mesh is adapted at every time step, the load of each process changes
as the shocks move. samurai redistributes the cells, and the fields defined on
them, with a load balancing module made of a driver and interchangeable
partitioning strategies:
\begin{itemize}
\item space-filling curves, with the Morton or the Hilbert ordering
  \cite{bader_space-filling_2013}. The cells are sorted along the curve, which
  is cut into $P$ segments of equal weight. The Hilbert curve gives more compact
  subdomains, the Morton keys are cheaper to compute;
\item diffusion. A discrete heat equation on the graph of the MPI neighbours
  \cite{cybenko_dynamic_1989} gives the load that each process sends to each
  of its neighbours, and the process cedes layers of cells along the
  corresponding interfaces. Cells only move between neighbours, and the balance
  is reached over several calls;
\item graph partitioning of the face-adjacent cells, which reduces the number of
  ghost cells, with ParMETIS \cite{karypis_parallel_1996} or PT-Scotch
  \cite{chevalier_pt-scotch_2008}. The adaptive repartitioning of ParMETIS
  starts from the current partition to limit the migration, while PT-Scotch
  partitions from scratch at every call;
\item a void strategy that moves no cell. It measures the fixed cost of the
  driver and gives the unbalanced reference run.
\end{itemize}
Each cell carries a weight, uniform by default, set per level or read from a
field. The driver rebalances when the imbalance $\max_p L_p / \bar{L} - 1$ of
the weighted loads $L_p$ exceeds a threshold, 5\% by default.

The double Mach reflection puts these strategies under conditions they handle
differently. The imbalance grows continuously as the shocks cross the domain,
the refined region is made of thin curved bands, and the jet concentrates cells
near the wall. At fixed mesh parameters and fixed $P$, we run the benchmark with
each strategy and measure the imbalance before and after each call, the number
of migrated cells, the time spent in partitioning and migration, and the effect
of the shape of the subdomains on the ghost updates and on the time to
solution. The rebalancing frequency is varied as well. Repeating the
comparison along the strong scaling and scaled-resolution series shows how the
balance between the quality of a partition, its cost and the data it moves changes
with $P$.

\paragraph{Strong scaling.}
The levels $(\ell_{\min}, \ell_{\max})$ and the threshold are fixed, and the
number of MPI processes $P$ increases from a reference count $P_0$. We report
the speedup and the parallel efficiency
\begin{equation}\label{eq:specs:samurai:dmr:strong}
  S(P) = \frac{T(P_0)}{T(P)}, \qquad E_s(P) = \frac{P_0}{P}\,S(P),
\end{equation}
where $T(P)$ is the time to solution on $P$ processes.

\paragraph{Scaled-resolution series.}
This series increases $\ell_{\max}$ by one each time $P$ is multiplied by 4.
Each refinement of the finest level multiplies the number of cells of the
uniform mesh by 4 in two dimensions, so a uniform mesh would keep the same
number of cells per process. On an adapted mesh, however, the number of cells
does not grow by a factor 4 per level. The refined cells follow
discontinuities, which are curves in the plane, and their share of the mesh
drops as the resolution increases. The number of time steps also doubles with
each level, through the CFL condition. The work per process therefore changes
along the series, so the series is not a weak scaling, and the time to
solution alone is not a fair measure. We report instead the relative throughput
per process
\begin{equation}\label{eq:specs:samurai:dmr:throughput}
  R(P) = \frac{N_{\mathrm{cu}}(P) / \left(P\, T(P)\right)}{N_{\mathrm{cu}}(P_0) / \left(P_0\, T(P_0)\right)},
\end{equation}
where $N_{\mathrm{cu}}(P)$ is the number of cell updates of the run, one cell
advanced by one time step counting as one update. This ratio normalizes by the
work actually done, which suits an adaptive mesh.

The number of cells at the initial and final times does not describe the work
accumulated over the run. To show how far each run is from constant work per
process, we also report the cell updates per process $N_{\mathrm{cu}}(P)/P$ and
the distribution of the load over the processes during the run, given by the
minimum, mean and maximum of the loads $L_p$ at each time step.

%%%%%%%%%%%%%%%%%%%%%%%%%%%%%%%%%%%%%%%%%%%%%%%%%%%%%%%%%%%%%%%%%%%%%%%%%%%%%%%%%%%%%%%%%%%%%%%%%%%%%%%%%%%%%%%%%%%%%%%%


\subsection{Benchmarking tools used}

The benchmark is the \texttt{euler\_2d} program of samurai-euler, which builds
samurai with MPI. The test case is selected with
\texttt{-{}-test-case double\_mach\_reflection}. The default final time of the
program is $0.25$, so the runs must set \texttt{-{}-Tf 0.2}. A run reads
\begin{verbatim}
mpiexec -n 64 ./euler_2d --test-case double_mach_reflection --Tf 0.2 \
        --order 2 --scheme hllc --min-level <lmin> --max-level <lmax> \
        --load-balancing-at <n> --metrics-file dmr_np64.json --timers
\end{verbatim}

With \texttt{-{}-load-balancing-at N}, the mesh adaptation rebalances the load
every $N$ adaptations, with the Hilbert curve and a uniform weight. The option
is 0 by default, which disables load balancing. The other strategies, the
weights and the per-call statistics of the driver (imbalance before and after,
migrated cells) are currently only available through the C++ API of samurai.
For the runs of this benchmark, we will make the choice of strategy available
on the command line of \texttt{euler\_2d}.

The option \texttt{-{}-metrics-file} writes the performance metrics of the run in
JSON. Each rank accumulates them on its own subdomain, and they are reduced
once at the end of the run: cell counts and cell updates are summed over the
ranks, and times are taken on the slowest rank. The file contains
\begin{itemize}
\item the time to solution, which covers the time loop only (initialization and
  file writing are measured but kept out of it);
\item the share of the time to solution spent in mesh adaptation;
\item the number of cell updates and the throughput in millions of cell updates
  per second (Mcu/s);
\item the number of cells and the sparsity index (cells of the adapted mesh over
  cells of the uniform mesh at $\ell_{\max}$) at the initial and final times.
\end{itemize}
The \texttt{-{}-timers} option of samurai adds a breakdown of the time spent in
mesh adaptation, ghost updates and the numerical scheme. It also reports the
time spent in load balancing, split between partitioning and migration, and,
for ParMETIS and PT-Scotch, the construction of the cell graph.

The regression tests of samurai-euler compare the double Mach reflection on a
uniform mesh with reference fields for the HLL, HLLC and Rusanov solvers, which
gives the \texttt{regression-test-pass} metric.


%%%%%%%%%%%%%%%%%%%%%%%%%%%%%%%%%%%%%%%%%%%%%%%%%%%%%%%%%%%%%%%%%%%%%%%%%%%%%%%%%%%%%%%%%%%%%%%%%%%%%%%%%%%%%%%%%%%%%%%%


\subsection{Input/Output Dataset Description}


\subsubsection{Input Data:}
  \begin{itemize}
  \item No external dataset: the domain, the initial state and the boundary
    conditions are defined in \texttt{euler/init/double\_mach.hpp}, and the
    mesh is generated from the command-line parameters.
  \item Parameters: mesh levels (\texttt{-{}-min-level}, \texttt{-{}-max-level}),
    multiresolution threshold (\texttt{-{}-mr-eps}), final time (\texttt{-{}-Tf}),
    Riemann solver (\texttt{-{}-scheme}), order (\texttt{-{}-order}), slope limiter
    (\texttt{-{}-slope-limiter}), CFL number (\texttt{-{}-cfl}), rebalancing
    frequency (\texttt{-{}-load-balancing-at}).
  \item Software: a tagged samurai release. The release used for the runs
    will be reported with the results.
  \end{itemize}

\subsubsection{Output Data:}
  \begin{itemize}
  \item The JSON metrics file and the timers report of each run.
  \item Optionally, the density, the pressure and the level of each cell in
    HDF5/XDMF format, to compare the solution with the published reference.
  \end{itemize}


%%%%%%%%%%%%%%%%%%%%%%%%%%%%%%%%%%%%%%%%%%%%%%%%%%%%%%%%%%%%%%%%%%%%%%%%%%%%%%%%%%%%%%%%%%%%%%%%%%%%%%%%%%%%%%%%%%%%%%%%


\subsection{Results summary}

No scaling result is reported in this version of the deliverable. The next
version will give, for each run, the machine (one of those listed in Appendix A),
the compiler and MPI versions, the number of processes per node, the load
balancing settings and the number of repetitions. It will show
\begin{itemize}
\item the density and the level of the cells at $T_f$;
\item for the strong scaling, the time to solution, the share of mesh
  adaptation, and $S(P)$ and $E_s(P)$ of \eqref{eq:specs:samurai:dmr:strong};
\item for the scaled-resolution series, the cell updates per process
  $N_{\mathrm{cu}}(P)/P$, the distribution of the load over the processes
  during the run, and the relative throughput per process $R(P)$ of
  \eqref{eq:specs:samurai:dmr:throughput};
\item for each load balancing strategy, the imbalance before and after
  rebalancing, the number of migrated cells, the time spent in partitioning and
  migration, and the time to solution.
\end{itemize}


%%%%%%%%%%%%%%%%%%%%%%%%%%%%%%%%%%%%%%%%%%%%%%%%%%%%%%%%%%%%%%%%%%%%%%%%%%%%%%%%%%%%%%%%%%%%%%%%%%%%%%%%%%%%%%%%%%%%%%%%


\subsection{Challenges identified}

\begin{itemize}
\item The refined cells move with the shocks, so a partition that is balanced
  at the start of the run does not stay balanced. The strong scaling depends on
  how often, and at what cost, the load is redistributed between the processes.
  The strategies trade the quality of the partition against its cost and the
  number of migrated cells, and the best trade-off may change with $P$.
\item The mesh is adapted at every time step. Its share of the time to
  solution is a direct measure of the overhead of the adaptation, and it must
  stay bounded as $P$ grows for the adaptive run to keep its advantage over a
  uniform mesh.
\item At fixed threshold, the sparsity index decreases with $\ell_{\max}$, so
  the scaled-resolution series cannot keep the work per process constant. The
  relative throughput per process of \eqref{eq:specs:samurai:dmr:throughput}
  normalizes by the work done, but the comparison is only fair if the cost of
  a cell update does not depend on the resolution. Reporting the load distribution
  during the run also requires the metrics file to record the per-process load
  over time, since it currently gives cell counts at the initial and final
  times only.
\end{itemize}
